%%%%%%%%%%%%%%%%%%%%%%%%%%%%%%%%%%%%%%%%%%%%%%%%%%%
%
%  New template code for TAMU Theses and Dissertations starting Spring 2021.  
%
%
%  Author: Thesis Office
%  
%  Last Updated: 1/13/2021
%
%%%%%%%%%%%%%%%%%%%%%%%%%%%%%%%%%%%%%%%%%%%%%%%%%%%
%%%%%%%%%%%%%%%%%%%%%%%%%%%%%%%%%%%%%%%%%%%%%%%%%%%%%%%%%%%%%%%%%%%%%
%%                           ABSTRACT 
%%%%%%%%%%%%%%%%%%%%%%%%%%%%%%%%%%%%%%%%%%%%%%%%%%%%%%%%%%%%%%%%%%%%%

\chapter*{ABSTRACT}
\addcontentsline{toc}{chapter}{ABSTRACT} % Needs to be set to part, so the TOC doesn't add 'CHAPTER ' prefix in the TOC.

\pagestyle{plain} % No headers, just page numbers
\pagenumbering{roman} % Roman numerals
\setcounter{page}{2}

Understanding human intent in natural language interactions remains a fundamental challenge in artificial intelligence, shaping the effectiveness, reliability, and safety of conversational systems. This research develops a unified approach to intent understanding that spans everyday cooperative communication, conversational failures, and adversarial intentions. We start with dialog acts, the minimal units of communicative intent, and introduce a multimodal framework for dialog act classification that integrates acoustic and linguistic cues to understand speaker intentions in real-time spoken conversations. Building on this foundation, we propose a contextual dialogue breakdown detection system designed for deployment in high-stakes, real-world conversational environments, enabling timely recovery when intent recognition fails. Crucially, intent is not always benign and the AI systems must also be able to recognize adversarial intents. We conduct a comprehensive evaluation of large language models for deception detection in multimodal settings, examining how AI systems can identify malicious intent across diverse domains including courtroom trials, interpersonal interactions, and online reviews. Finally, this exploration culminates in a novel jailbreak defense for LLMs that treats adversarial attacks as a form of intent manipulation. This proposed approach incorporates human-inspired deception resistance strategies through a combination of supervised fine-tuning and reinforcement learning. Taken together, this work presents a unified view of intent, from cooperative goals to deceptive and adversarial behavior, showing how advances in foundational dialog understanding can directly inform robustness and safety strategies, and enabling conversational AI that is more accurate, trustworthy, and resilient.
\pagebreak{}
